\BourbakiUnnumberedChapter{历史注记}{历史注记（第四、五章）}

（第四章和第五章）

（方括号中的数字指本注记末尾的参考文献。）

交换域的理论——以及与之密切相关的多项式理论——直接源于直到19世纪中叶仍构成经典代数学主要对象的内容：代数方程的求解，以及与之等价的几何作图问题。

一旦人们试图求解一个次数为 \(>1\) 的代数方程，就会发现自己面临着全新的计算困难，因为不再可能通过对数据进行“有理”运算来确定未知量。这一困难必定在很早的阶段就被察觉到了；巴比伦人最重要的贡献之一，应当是他们能够将二次和四次方程的求解归结为一种新的代数运算，即开平方（正如流传下来的文本中许多以这种方式求解的数值方程所表明的那样（{[}1{]}，第183—193页））。就形式计算而言，古代在代数方程求解问题上从未超越这一点；古典时期的希腊人实际上仅限于用几何语言重述巴比伦公式，而它们以代数形式的使用在希罗（公元100年）和丢番图之前没有文献记载。

希腊人是在一个相当不同的方向上取得了决定性的进展。关于巴比伦人如何构想和计算非完全平方整数的平方根，我们所知甚少*；在流传下来的关于这一问题的罕见文本中，他们似乎满足于相当粗略的近似方法（{[}1{]}，第33—38页）。毕达哥拉斯学派曾以相当严格的方式确定了可公度量的概念，并赋予其近乎宗教的性质，却未能始终坚持这一观点；也许正是反复尝试有理地表示 \(\sqrt{2}\) 的失败，最终导致他们证明了该数是无理数**。

我们已在别处说过（参见第三卷第四章的历史注记），这一发现如何成为数学史上的分水岭，深刻地影响了希腊人对“数”的概念，并促使他们创造了一种纯粹几何特征的代数，以便为不可公度比找到一种表示方式（或者也许是一种“存在性”证明），因为他们拒绝将这些比视为数。在大多数情形，他们把一个代数问题归结为两条适当选取的辅助平面曲线的交点，或归结为相继若干次确定这样的交点。一些权威性甚低的晚期传统把这类作图的第一个分类的引入归功于柏拉图，这一分类注定有着漫长而辉煌的历程：似乎出于哲学而非数学的原因，所谓的“尺规作图”，即其中出现的仅有的辅助曲线为直线和圆的那些作图，被单独列为一类*。无论如何，欧几里得在他的《几何原本》 {[}2{]} 中仅限于专门处理可以这种方式解决的问题（尽管没有用特定名称来刻画它们）；这一情况无疑极大地促使后世数学家将注意力集中在这些问题上。但我们今天知道**，可以用“尺规”求解的代数方程是一种非常特殊的类型；特别地，一个不可约方程

立即给出 \(\sqrt{5}\) 的无理性的一个证明，而且他提出了一个假设（遗憾的是没有任何文本支持这一假设），即毕达哥拉斯学派正是以这种方式发现了无理数（K. von FRITZ，Ann. of. Math.，第46卷（1945年），第242页）。

* 关于这一原则，人们还将平面曲线分类为“平面轨迹”（直线和圆）、“立体轨迹”（圆锥曲线，通过平面截取立体——圆锥得到）的作法归功于柏拉图，所有其他曲线则统称为“τόποι γραμμικοι”。值得注意的是，这一分类的影响甚至延伸到了笛卡尔，他在其《几何学》中将次数为2n−1和2n的方程归入同一“属”，无疑是因为次数为1或2的方程通过“平面轨迹”的交点求解，而次数为3或4的方程通过“立体轨迹”的交点求解。

** 确定一条直线与一个圆（或两个圆）的交点，等价于求一个二次方程的解，该方程的系数是所考虑的直线和圆（或两个圆）的方程系数的有理函数。由此容易推出，从给定点“用尺规”作出的点的坐标属于有理数域Q的一个扩域L，其构造如下：若K是通过将给定点的坐标添加到Q而得到的域，则存在一个递增序列 \((\mathrm{L}_{i})_{0\leq i\leq n}\) ，其中的域都介于K和L之间，满足条件 \(K=L_{0}\) 、\(L=L_{n}\) 、\([L_{i}:L_{i-1}]=2\) （\(1\leq i\leq n\) ）。通过对n进行归纳，我们推出由L生成的伽罗瓦扩域N在K上的次数是2的幂；反之可以证明，若此条件成立，则存在一个序列 \((\mathrm{L}_{i})\) ，其中各域介于K和L之间且具有上述性质，因此所提出的问题可以用尺规求解（参见N. TSCHEBOTARÖW，Grundzüge der Galoisschen Theorie（H. Schwerdtfeger译），格罗宁根（P. Noordhoff出版社），1950年，第351页）。

三次（在有理数域上）的方程不能用这种方式求解，而希腊人很早就遇到了一些至今仍很著名的三次问题，例如倍立方（求解 \(x^3 = 2\) ）以及角的三等分；另一方面，化圆为方问题使他们面对一个超越性问题。我们发现，为了解决这些问题，他们引入了许多曲线，既有代数曲线（圆锥曲线、狄奥克勒斯的蔓叶线、尼科梅德斯的蚌线），也有超越曲线（狄诺斯特拉托斯的割圆曲线、阿基米德螺线）；所有这些不可避免地促使他们对这些不同的曲线进行独立研究，从而为解析几何、代数几何以及无穷小演算的未来发展铺平了道路。但这些方法对代数方程的求解没有带来任何进展*，而古代唯一对这一问题作出显著贡献并对中世纪和文艺复兴时期的代数学家产生持久影响的著作，是欧几里得《几何原本》 {[}2{]} 的第十卷；在这本书中（其中主要成果被某些历史学家归于泰阿泰德），他考察了由多个根式组合得到的表达式，例如 \(\sqrt{\sqrt{a} \pm \sqrt{b}}\) （\(a\) 和 \(b\) 为有理数），给出了这些表达式为无理数的条件，将它们分成众多的类别（并证明这些类别互不相同），并研究了这些不同的无理数之间的代数关系，例如我们今天会写成

\[
\sqrt {\sqrt {p} + \sqrt {q}} = \sqrt {\frac {1}{2} (\sqrt {p} + \sqrt {p - q})} + \sqrt {\frac {1}{2} (\sqrt {p} - \sqrt {p - q})};
\]

所有这些都用《几何原本》中惯用的几何语言表述，这使得论述显得尤为繁密而艰涩。

古典希腊数学衰落之后，与代数方程相关的概念经历了演变。毫无疑问，在整个古典时期，希腊人掌握了平方根无限逼近的方法，但遗憾的是，我们对这些方法知之甚少**。

在印度人那里，后来又在阿拉伯人及其中世纪的西方效仿者那里，各阶开方成为一种运算，它倾向于被看成与代数的有理运算具有同等地位的基本运算，并且像后者一样，用越来越便于在计算中处理的符号来表示*。二次方程的理论在中世纪得到了完善（根的个数、负根、不可能情形、重根），四次方程的理论也是如此，它们为“用根式”求解方程提供了公式的范例，几个世纪中代数学家们都试图仿照这些范例，为更高次方程的求解建立类似的公式，首先便是三次方程。比萨的莱昂纳多在13世纪主要负责将阿拉伯科学引入西方，他认识到，无论如何，欧几里得在第十卷中分类的无理数不能用于此目的（在一个充满这类证明的理论中，这又是一个新的不可能性证明），并且我们看到他已经对立方根尝试了类似的计算，得到了诸如以下的关系：

\[
\sqrt [ 3 ]{1 6} + \sqrt [ 3 ]{5 4} = \sqrt [ 3 ]{2 5 0},
\]

类似于欧几里得关于平方根的公式（我们实际上在阿拉伯人那里能找到更早的例子）。但还要经过三个世纪徒劳无功的努力，直到16世纪初，希皮奥内·德尔·费罗才最终得到方程 \(x^3 + ax = b\) 的求解公式：

\[
x = \sqrt [ 3 ]{\frac {b}{2} + \sqrt {\left(\frac {b}{2}\right) ^ {2} + \left(\frac {a}{3}\right) ^ {3}}} + \sqrt [ 3 ]{\frac {b}{2} - \sqrt {\left(\frac {b}{2}\right) ^ {2} + \left(\frac {a}{3}\right) ^ {3}}}.\tag{1}
\]

我们在此不描述这一轰动性发现中引人入胜的一面——它引发的塔尔塔利亚一方与卡尔达诺及其学派之间的争执——也不描述作为主角的学者们那些往往饶有趣味的人物形象。但我们必须指出卡尔达诺及其弟子们在方程论中带来的决定性进展。卡尔达诺对使用负数的反感不如其同时代的大多数人强烈，他因此观察到三次方程可能有三个根，而四次方程有四个根（{[}3{]} ，第IV卷，第259页），并指出 \(x^3 + bx = ax^2 + c\) （他知道如何消去其中 \(x^2\) 项的一个方程）的三个根之和总等于 \(a\) （同上）。无疑受这一关系及其一般特征的直觉引导，他最先有了根的重数这一概念的初步想法；尤其是，他大胆地（并非没有修辞上的谨慎）对含有负数平方根的表达式进行形式计算。他很可能是被下述事实引到这一步的：当 \(\left(\frac{b}{2}\right)^{2} + \left(\frac{a}{3}\right)^{3} < 0\) 时（所谓的“不可约情形”，卡尔达诺在其中认识到三个实根的存在），这类表达式在使用公式(1)时自然出现；无论如何，这一点由他的弟子R. 邦贝利清楚地表明，他在其《代数学》（{[}4{]} ，第293页）中证明了关系

\[
\sqrt [ 3 ]{2 + \sqrt {- 1 2 1}} = 2 + \sqrt {- 1}
\]

并且他注意以已经非常接近现代阐述的形式明确给出复数运算的规则*。最后在1545年，卡尔达诺的另一个学生L. 费拉里借助一个三次的辅助方程**成功地求解了一般的四次方程。

在如此迅速的进展之后，接下来的时期，直到18世纪中叶，几乎未曾发展意大利学派引入的新思想。得益于他在代数记号方面带来的根本性进步，韦达得以一般性地表达代数方程的系数与根之间的关系，至少当根全为正数时如此***（{[}5{]} ，第158页）。更大胆的A. Girard {[}6{]} 毫不犹豫地断言（当然没有证明）：次数为 \(n\) 的方程恰好有 \(n\) 个根，只要把那些“不可能的根”计入，每个根按其重数计算，并且这些根满足韦达给出的关系；他还第一次得到了根的等幂和的表达式，直至指数4。

但17世纪的精神转向了其他方向，代数只是间接地从解析几何和无穷小演算的新发现中获益。因此，笛卡尔获得

* 邦贝利（{[}4{]}，第169页和第190页）将复数视为四个基元素的具有正系数的“线性组合”：“piu”（+1）、“meno”（-1）、“piu de meno”（+i）和“meno de meno”（-i）；特别地，他提出一条公理，即“piu”与“piu de meno”不能相加，这是线性无关概念的首次出现。

** 方程被化简为形式 \(x^4 = ax^2 + bx + c\) ，我们确定一个数 \(z\) 使得方程的右边

\[
(x ^ {2} + z) ^ {2} = (a + 2 z) x ^ {2} + b x + (c + z ^ {2})
\]

是完全平方，这就给出了关于 \(z\) 的一个三次方程。

韦达是古人的热情仰慕者，他系统性地避免在论证中引入负数；尽管如此，他有时仍能用自己的语言表达系数与根之间的关系，即使某些根为负；例如，若方程 \(x^3 + b = ax\) 有两个正根 \(x_1, x_2\) （ \(a > 0, b > 0\) ），韦达表明 \(x_1^2 + x_2^2 + x_1x_2 = a\) 与 \(x_1x_2(x_1 + x_2) = b\) （{[}5{]} ，第106页）。

代数曲线的切线（参见《历史注记》第四卷，第一至三章，第46页）与他的弟子胡德所陈述的代数方程根的重数判别准则有关（{[}7{]}，第433页及第507-509页）。无疑，代数函数与超越函数之间的区分也应归功于笛卡尔的影响，这与他在《几何学》中所引入的“几何”曲线与“机械”曲线之间的区分相平行（参见《历史注记》第四卷，第一至三章，第46页及第61页）。无论如何，这一区分在J.格雷戈里那里已表述得十分清楚，他在1667年甚至试图证明圆扇形面积不可能是弦和半径的代数函数*。“超越”这一术语出自莱布尼茨，他对这些分类问题在整个学术生涯中始终兴趣不减，并且大约在1682年，他通过证明sin x不是x的代数函数，发现了格雷戈里所追求结果的一个简单证明（{[}8{]}，第五卷，第97-98页）**。莱布尼茨与他的朋友契恩豪斯一道，实际上是当时少数仍对代数方程“根式”求解问题感兴趣的数学家之一。在其学术生涯初期，我们发现他研究三次方程的“不可约情形”，并使自己相信（事实上，其证明并不充分）在这种情况下不可能从解的公式中消去虚量（{[}9{]}，第547-564页）。大约在同一时期，他还尝试用根式求解五次方程，但同样未获成功；后来当契恩豪斯声称通过形如y = P(x)的变换（其中P是适当选取的四次多项式）使方程除首尾两项外的所有项都消失，从而解决该问题时，莱布尼茨立刻注意到确定P(x)系数的方程次数为 \textgreater{} 5，并认为该方法注定失败（{[}9{]}，第402-403页）。

似乎是新分析的需求逐渐重新点燃了对代数的兴趣。有理分式的积分（由莱布尼茨和约翰·伯努利完成）以及与之密切相关的虚对数问题，为更深入地发展虚数运算提供了机会，并重新提起把多项式分解为一次因式的问题（“代数基本定理”）*。就在18世纪初，二项式方程 \(x^n - 1 = 0\) 的求解被科茨和棣莫弗化为将圆分成 \(n\) 等份的问题；因此，为了得到其根的“根式”表达式，只需知道当 \(n\) 为奇素数时如何做到这一点，而棣莫弗指出，代换 \(y = x + \frac{1}{x}\) 把问题归结为求解次数为 \((n - 1)/2\) 的一个方程的“根式”解。至于“基本定理”，在一般“根式”解的多次失败尝试（包括欧拉的几次尝试 {[}10{]} ）之后，人们开始倾向于先验的证明，而不使用解的显式公式。这里不深入讨论所提出方法的细节（这些方法后来在拉格朗日和高斯的证明中结出果实；参见第二卷第六、七章及第三卷第八章的历史注记），只需在此指出18世纪中叶人们看待这一问题的视角；人们承认（除了对一般性的模糊感觉之外没有任何别的理由，这种感觉无疑同A. 吉拉尔那里一样，来自系数与根之间存在关系的事实）次数为 \(n\) 的方程总有 \(n\) 个“理想的”根，对它们可以像对数那样进行运算，而无需知道它们是否是（实数或复）数；需要证明的是（必要时利用对理想根的运算），这些根中至少有一个是通常的复数**。在这种有缺陷的形式中，可以辨认出“形式添加”这一一般观念的最初萌芽，尽管有高斯的反对（{[}13{]} ，第III卷，第1页），它后来成为现代交换域理论的基础。

随着拉格朗日（{[}11{]} ， \(a\) ）与范德蒙德 {[}12{]} 的基本回忆录的发表，1770年开启了代数方程历史上一个新的决定性时期。此前一直独占统治地位的、多少带有偶然性的寻找求解公式的经验主义，被对所提问题以及能够解决它们的方法的系统分析所取代；这一分析在六十年后导致了伽罗瓦的最终成果。拉格朗日和范德蒙德都是从次数 \(\leqslant 4\) 的方程求解公式中根式的多重取值所引入的歧义出发的；这一事实早已引起欧拉（{[}10{]} ， \(a\) ）的注意，他除其他事情外还证明了，在德尔·费罗公式中应当如何把其中出现的根式关联起来，以便得到3个根而不是9个。拉格朗日指出，德尔·费罗公式中的每个立方根都可以写成形式 \(\frac{1}{3} (x_1 + \omega x_2 + \omega^2 x_3)\) ，其中 \(\omega\) 是一个单位立方根，而 \(x_1, x_2, x_3\) 是所给方程的根按某种顺序的排列，他还作出了一个基本的观察：三个根的函数 \((x_1 + \omega x_2 + \omega^2 x_3)^3\) 在这三个根的每一置换下只能取两个不同的值，这先验地解释了该方程的求解方法之所以成功。对四次方程求解方法的类似分析把他引导到四个根的函数 \(x_1 x_2 + x_3 x_4\) ，该函数在根的任意置换下仅取三个不同的值，因而是某个三次方程的根，其系数是给定方程系数的有理函数*；正如拉格朗日所说，这些事实构成了“解三次和四次方程的真正原理，也可以说是其形而上学**”（{[}11{]} ，a)，第357页）。在这些例子的引导下，他计划就一般情形研究：对于次数为 \(n\) 的方程，当根被任意置换时，根的有理函数 V 可能取到的值的个数 \(\nu\) ***（V 是这 \(n\) 个根的有理函数）；他由此实际上开创了（其术语仍紧密贴合方程论）群论和域论，并且已经运用与今天所用的相同原理得到了若干基本结果。例如，他证明了数 \(\nu\) 整除 \(n!\) ，所用的推理正是今天用来证明有限群的子群的阶整除该群的阶的那一推理。更值得注意的是他证明的定理：若 \(V_1\) 与 \(V_2\) 是根的两个有理函数，使得 \(V_1\) 与 \(V_2\) 在同样的诸置换下保持不变，则其中每一个都是另一个与方程系数的有理函数（这是伽罗瓦定理的一个特例，该定理把伽罗瓦扩张的子扩张刻画为其伽罗瓦群的不变量域）：“这个问题，”他说，“在我看来是方程论中最重要的问题之一，我们将给出的一般解将有助于为代数的这一部分带来新的光明。”（{[}11{]} ，a)，第374页）。

所有这些研究在拉格朗日看来自然只是预备性工作，为的是分析通过逐次化归为较低次的方程来求解代数方程的各种可能方法；正如他所证明的，这样的方法关系到构造根的有理函数，使其在根的置换下取少于 \(n\) 个值。无疑受到三次方程结果的引导，他就一般情形引入了“拉格朗日预解式” \(y_{k} = \sum_{h=1}^{n} \omega_{k}^{h} x_{h}\) ，其中 \(\omega_{k}\) 是 n 次单位根 \((1 \leqslant k \leqslant n)\) ，并清楚地表明这 n 个数的知识如何蕴含根 \(x_{k}\) 的知识，且一般地考察了 \(y_{k}\) 所满足的方程的次数；例如他证明，若 n 是素数，则诸 \(y_{k}^{n}\) 是一个 n-1 次方程的根，其系数是某个方程的一个根的有理函数，该方程的次数为 \((n-2)!\) ，而这些系数又可由所给方程的系数有理地表示。“这里，除非我弄错，”他总结道，“是方程求解的真正原理，也是最适于引导我们达到目的的分析；正如我们所见，一切都归结为一种组合演算，由此可以先验地求得应当期待的结果”（{[}11{]} ，a)，第403页）。

关于范德蒙德的回忆录，尽管独立于拉格朗日的著作，却在许多点上与后者有所接触，特别是在研究根的有理函数、使其在根的置换下取尽可能少的不同值这一思想*，以及他为此目的同样引入的“拉格朗日预解式”的研究上。他的工作远不具有拉格朗日工作的清晰性和一般性；然而在一点上他明显走得更远，即把同样的思想应用于分圆方程 \(x^n - 1 = 0\) ，其中 \(n\) 为奇素数。拉格朗日满足于指出这个方程可化为次数为 \(m = (n - 1)/2\) 的有理系数方程，而不试图在 \(n \geqslant 11\) 时求解它；范德蒙德则断言，后一方程的各拉格朗日预解式的 \(m\) 次幂是有理的，这是由 \(x^n - 1 = 0\) 的各个根之间的关系得出的；但他仅限于对 \(n = 11\) 验证这一断言的可靠性，而不作一般的论证。

直到30年后，范德蒙德宣布的结果才由C. F. 高斯**完全证明。他关于方程 \(x^n - 1 = 0\) （ \(n\) 为奇素数）的决定性成果，属于他那令人难忘的算术研究的总体计划（{[}13{]} ，第I卷，第413页及以下），它们尤其展示了他驾驭我们现今所称的循环群理论的娴熟技巧。在证明了多项式 \(\Phi_n(x) = (x^n - 1)/(x - 1)\) 对奇素数 \(n\) 不可约***之后，他想到把它的 \(n - 1\) 个根写成形式 \(\zeta^{g^{k}} = \zeta_{k} \quad (0 \leqslant k \leqslant n - 2)\) ，其中 \(g\) 是同余式 \(z^{n-1} \equiv 1 \pmod{n}\) 的一个本原根（用现代术语说，这相当于揭示出群 \(\Gamma\) ——方程 \(\Phi_{n}(x) = 0\) 的群——是循环群这一事实）。对每个除数 \(e\) （它整除 \(n - 1\) ），他关联起 \(f = (n - 1)/e\) 个“周期” \(\eta_{\nu} = \zeta_{\nu} + \zeta_{\nu + e} + \zeta_{\nu + 2e} + \cdots + \zeta_{\nu + (f - 1)e}\) （ \(1 \leqslant \nu \leqslant f\) ），并且他本质上证明了诸 \(\eta_{\nu}\) 的具有有理系数的线性组合构成一个域，它由 \(f\) 个周期 \(\eta_{\nu}\) 中的任何一个生成，且在有理数域上的次数为 \(f\) （这个域当然对应于 \(\Gamma\) 的阶为 \(e\) 的子群）。我们在此无法详述这一分析及其所蕴含的重要算术结论；只需指出，它特别使他得到了那个著名的定理，即边数等于形如 \(2^{2^{k}} + 1\) 的素数的正多边形可以用直尺和圆规作出*。至于方程 \(\Phi_{n}(x) = 0\) 的根式解，把周期理论应用于拉格朗日预解式的 \(f\) 次幂 \(\sum_{\nu=0}^{f-1} \omega^{\nu} \eta_{\nu}\) （其中 \(\omega^{f} = 1\) ）即可容易地得到**。

拉格朗日的同胞鲁菲尼的研究与《算术研究》同时代，直接承自拉格朗日；他们从拉格朗日停下的地方接着研究这一问题，目标在于证明“一般”***五次方程“用根式”求解的不可能性。鲁菲尼的证明冗长而晦涩，尽管多次修订，仍不完整；但它已经非常接近阿贝尔后来得到的（原则上正确的）证明。其主要意义在于引入了代换演算和群论的最初步概念，这

鲁菲尼发展了它们以表明，不存在方程5个根的一个函数，当根被任意置换时，它取多于2个且少于5个值。

我们已经在第一章的历史注记中说过，置换群理论的这第一个轮廓是如何在数年之后由柯西加以发展和系统化的。但是，尽管发展拉格朗日关于代换思想的必要概念由此逐渐明晰，仍然有必要以同样清晰的方式阐述域理论的基本原理。这正是鲁菲尼所缺乏的，而阿贝尔和伽罗瓦将在代数方程求解问题的最后阶段着手进行这项工作。

阿贝尔在其短暂的一生中始终未曾停止对这一问题的萦怀。几乎还在孩提时代，他就相信自己找到了一般五次方程根式求解的公式。后来，当他认识到自己的错误时，便不曾停歇，直到成功证明了这样的公式并不存在（{[}14{]} ，第I卷，第166页）。但他并未止步于此。当他的竞争者雅可比以分析学家的方式发展椭圆函数理论时，支配阿贝尔在这一问题上的工作的却是代数观点，其中心是椭圆函数分割方程的理论（{[}14{]} ，第I卷，第265、377页及多处）。他由此得到了一些新型的可根式求解的方程，其方法仿照高斯处理分圆方程的方法（{[}14{]} ，第I卷，第310和358页）*；从这一结果他上升到“阿贝尔”方程的概念，并在一篇著名的回忆录中证明了这类方程可根式求解（{[}14{]} ，第I卷，第478页）；正是在这个场合，他精确定义了给定域（由他所研究的方程的系数生成的域）上不可约多项式的概念**。最后，死神于1829年将他击倒，当时他正在攻克刻画所有可根式求解的方程这一一般问题，并且刚刚向克雷尔和勒让德通报了已非常接近伽罗瓦结果的研究（{[}14{]} ，第II卷，第219-243、269-270和279页）。

三年后，为这座大厦加冕的任务正是留给了他 \([15]\) 。与阿贝尔类似，但方式更为精确，他首先定义了（除名称外）一个量属于由给定量生成的域、添加的概念以及给定域上的不可约多项式。给定一个方程 \(F(x) = 0\) ，它没有重根，系数属于一个给定域 \(K\) ，他相继表明如何“总可以构造根的一个函数 \(V\) ，使得以一切可能的方式在该函数中置换根所得到的值没有一个等于另一个”，该函数“具有这样的性质：所给方程的所有根都可以表示为 \(V\) 的有理函数”，而且，若 \(V, V', V'', \ldots\) 是 \(V\) 所满足的不可约方程的根，则“若 \(a = f(V)\) 是所给方程的一个根，那么 \(f(V')\) 也是”（{[}15{]} ，第36-37页）；用现代术语说，他由此证明了 \(V\) ，同它在 \(K\) 上的任一共轭元一样，生成域 \(N\) （ \(F\) 的根的域）。然后他定义群 \(\Gamma\) （即 \(F\) 的群）为根 \(x_i\) 的通过代换 \(V\) 而得到的那些置换的集合：在把每个 \(x_i\) 表示成 \(V\) 的函数的有理表达式中，代入 \(V\) 的任何一个共轭元；他随即得到 \(K\) 的元素由在 \(\Gamma\) 的每一置换下不变这一性质所作的基本刻画（{[}15{]} ，第38-39页）。然后他证明，若 \(N\) 包含另一个多项式的根域 \(L\) ，则 \(N\) 在 \(L\) 上的群是 \(\Gamma\) 的一个正规子群（他在这个场合引入了这一概念）（{[}15{]} ，第41页及第25-26页）。由此他最终导出方程可用根式求解的判别准则，其推理在本质上如下：设基域 \(K\) 包含所有单位根，则由假设存在中间域的一个递增序列 \((K_i)_{0 \leq i \leq m}\) ，它们介于 \(K\) 与 \(N\) 之间，满足 \(K_0 = K\) 、\(K_m = N\) ，其中域 \(K_{i+1}\) 是通过向 \(K_i\) 添加一个二项方程 \(x^{n_i} - a_i = 0\) （ \(a_i \in K_i\) ）的所有根而得到的。于是在 \(\Gamma\) 中存在子群的一个递降序列 \((\Gamma_i)\) ，使得 \(\Gamma_0 = \Gamma\) 、\(\Gamma_m = \{\varepsilon\}\) （单位元），\(\Gamma_{i+1}\) 在 \(\Gamma_i\) 中正规，且商群 \(\Gamma_i / \Gamma_{i+1}\) 是循环群（此时称群 \(\Gamma\) 为可解的）。反之，当这一条件成立时，利用拉格朗日预解式可以表明 \(K_{i+1}\) 是由向 \(K_i\) 添加一个二项方程的所有根而得到的，因此方程 \(F(x) = 0\) 可用根式求解*。次数为 n \textgreater{} 4 的“一般”方程无法用根式求解这一不可能性，乃是该方程的群 \(\Gamma\) ——它同构于对称群 \(S_n\) （V，第59页，例5）——不可解这一事实的推论（I，第142页，习题10及第143页，习题16）。

自19世纪中叶起，代数学家们，正如我们已指出的（参见第一章的历史注记），大大扩展了他们的研究领域，此前该领域几乎完全局限于方程的研究。在伽罗瓦发现的启示下，人们认识到“用根式”求解的问题只是对无理数进行分类这一普遍问题的一个特殊且略显人为的情形。正是后者，在整个19世纪最后几十年里，从各个方面受到探讨，众多零散的结果逐渐积累，为施泰尼茨的综合工作铺平了道路。

首先，就代数无理数而言，分类的基本原则由伽罗瓦理论提供，该理论把代数方程的研究归结为其群的研究。因此，首先是置换群理论——我们不打算在此论述（参见第一章历史注记）——在纯代数中构成了该时期研究的主要对象。代数域理论中其余的进展，则归因于数论与代数几何的同时发展。实际上，这些进展主要涉及理论的阐述，且大多应归功于戴德金 {[}18{]} ，他引入了域和环的概念*，并且（结合他关于超复数系的研究）系统发展了扩张理论的线性方面（{[}18{]} ，第3卷，第33页及以下）。也是他，把伽罗瓦群看成所考虑的扩张的自同构群，而不仅仅是方程根的置换群，并且他（对数域）证明了关于自同构线性无关的基本定理（{[}18{]} ，第3卷，第29页）以及伽罗瓦扩张正规基的存在性（{[}18{]} ，第2卷，第433页）。最后，他着手处理无限次代数扩张的问题，并指出伽罗瓦理论不能照原样应用（伽罗瓦群的任意子群未必总等于扩张关于某个子扩张的群），而且凭借大胆的直觉，他已经尝试把伽罗瓦群看作一个拓扑群**——这一思想直到1928年才在克鲁尔发展的无限次伽罗瓦扩张理论中成熟（{[}24{]} ）。

与这一发展并行，域上超越元的概念被精确化。超越数的存在性第一次由刘维尔在1844年证明，他借助一个基于丢番图逼近理论的显式构造程序完成了证明 {[}17{]} 。康托尔在1874年利用关于集合的势的简单思考给出了另一个“非构造性”证明（参见V，第147页，习题3）；最后，埃尔米特在1873年证明了 e 的超越性，林德曼在1882年用类似于埃尔米特的方法证明了 \(\pi\) 的超越性，从而最终结束了关于古老的化圆为方问题的讨论*。

关于超越数在代数计算中的作用，克罗内克在1882年指出，若 x 在域 K 上是超越的，则域 \(\mathrm{K}(x)\) 同构于有理分式域 \(\mathrm{K}(\mathrm{X})\) （{[}19{]} ，a)，第7页）。他实际上把未定元向域的添加作为其代数数理论阐述的基石（{[}19{]} ，a)）。另一方面，戴德金和韦伯在同一年（{[}18{]} ，第I卷，第238页）表明了如何用算术方法建立代数曲线的理论。算术与代数几何之间的类比——它后来对双方都极为富有成效——就是这样从几个方向显现出来的。

在所有这些研究中，出现的域都由经典数学意义上的“具体”元素组成——（复）数或复变量的函数**。但早在1882年，克罗内克就完全意识到了这一事实（高斯和伽罗瓦只是隐约地瞥见）：在他的理论中，“未定元”只起代数的基元素的作用，而不是分析意义下的变量（{[}19{]} ，a)，第93-95页）；并且在1887年，他发展了这一思想，把它同一个宏大的纲领结合起来，该纲领的目标不外乎通过拒绝一切不能归结为整数上代数运算的东西，为整个数学重新奠基（参见第一册第四章的历史注记）。正是在这个场合，他接过了柯西 {[}16{]} 的一个想法——柯西曾把复数域 C 定义为剩余类域 \(\mathbf{R}[\mathbf{X}] / (\mathbf{X}^2 + 1)\) ——而克罗内克表明，代数数的理论如何完全独立于“代数基本定理”，甚至独立于实数理论，因为每个（有限次的）代数数域都同构于一个剩余类域 \(\mathbf{Q}[\mathbf{X}] / (f)\) （其中 \(f\) 是 \(\mathbf{Q}\) 上的一个不可约多项式）（{[}19{]} ，b)）。正如H. Weber {[}20{]} 几年后在发展域的公理化理论的第一个概要时所指出的，克罗内克的这个方法实际上适用于任何基域K。韦伯特别指出，可以把域 \(\mathbf{Z}/(p)\) （ \(p\) 为素数）取作K，从而使模 \(p\) 的同余演算被纳入域论之中；后者在18世纪下半叶经欧拉、拉格朗日、勒让德和高斯之手成形，它同代数方程理论所呈现的类比早已被人们注意到；在推进这一类比时，伽罗瓦（鉴于他在群论方面的研究）毫不犹豫地引入了“理想根”，即模 \(p\) 的一个不可约同余式的根*，并指出了其主要性质（{[}15{]} ，第15-23页）**。当把克罗内克的方法应用于 \(\mathbf{Z}/(p)\) 时，所得到的结果（除名称外）就是塞尔雷特和戴德金（{[}18{]} ，第I卷，第40页）已经对这些“伽罗瓦虚数”给出过的表述。

在世纪之交，除了所有这些“抽象域”的例子之外，又增添了韦罗内塞引入的形式幂级数域 {[}21{]} ，尤其是亨泽尔的 \(p\) 进数域 {[}22{]} 。正是后者的发现引导施泰尼茨（正如他明确说过的）把这些理论共同的概念分离出来，写进一篇基础性论文 {[}23{]} ，这篇论文可以被视为奠定了当今的代数学观念。系统地展开交换域公理的推论，他引入了素域、可分（代数）元、完全域等概念，定义了扩张的超越次数，并最终证明了任意域的代数闭扩张的存在。

更近的时候，施泰尼茨的理论已在几个重要方面得到完善。一方面，阿廷的工作凸显了伽罗瓦理论的线性特征 \([25]\) 。另一方面，导子的一般概念（模仿经典微分演算的形式性质）——戴德金已有预感（{[}18{]} ，第2卷，第412页），施泰尼茨在有理分式域的特殊情形下引入（{[}23{]} ，第209-212页）——已被成功地用于（对现代代数几何至关重要的）超越扩张的研究，特别是用于把可分性概念推广到超越扩张 \([26]\) 。

